\section{工程实践建议}

\subsection{开发高性能代码的工作流}

基于本课的内容，Hashimoto 教授推荐以下工作流：

\begin{enumerate}
    \item \textbf{先用 PyTorch 写正确的代码}：确保逻辑正确，编写测试用例
    \item \textbf{基准测试}：测量端到端性能，建立 baseline
    \item \textbf{Profile}：找到真正的瓶颈（不要凭直觉猜测）
    \item \textbf{尝试 torch.compile}：一行代码可能就够了
    \item \textbf{必要时写 Triton 内核}：仅对 Profile 确认的瓶颈操作
    \item \textbf{重新基准测试}：验证优化效果
\end{enumerate}

\begin{warningbox}{避免过早优化}
``有些同学花了三小时优化一个他们\textbf{以为}是瓶颈的组件，结果发现根本不是。'' 始终让 Profiling 数据指导你的优化方向。
\end{warningbox}

\subsection{CPU-GPU 异步执行的实践影响}

CPU-GPU 异步执行模型对日常开发有几个重要影响：

\begin{importantbox}{避免不必要的 CPU-GPU 同步}
以下操作会强制同步，应谨慎使用：
\begin{itemize}
    \item \texttt{print(tensor.item())} 或 \texttt{print(loss)}：需要将 GPU 数据传回 CPU
    \item \texttt{if tensor > threshold:}：需要在 CPU 上评估条件
    \item \texttt{tensor.numpy()}：GPU tensor 到 NumPy 数组的转换
    \item 任何在 CPU 上访问 GPU tensor 值的操作
\end{itemize}
\end{importantbox}

\subsection{矩阵维度选择}

来自上一节课的关键建议，在实践中极为重要：
\begin{itemize}
    \item 矩阵维度选择 \textbf{64/128/256 的倍数}
    \item 避免质数或不规则的维度
    \item nanoGPT 案例：词表大小从 50257 $\rightarrow$ 50304（最近的 64 的倍数），获得 25\% 加速
\end{itemize}

\subsection{选择正确的数值精度}

\begin{table}[H]
\centering
\begin{tabular}{lcc}
\toprule
\textbf{操作类型} & \textbf{推荐精度} & \textbf{原因} \\
\midrule
矩阵乘法 & FP16/BF16 输入，FP32 累加 & Tensor Core 原生支持 \\
Pointwise（ReLU, GeLU 等） & FP16/BF16 & 减少内存带宽需求 \\
归约（sum, softmax, norm） & FP32 累加 & 避免小值累加误差 \\
Loss 函数 & FP32 & 需要大动态范围 \\
\bottomrule
\end{tabular}
\caption{不同操作类型的推荐数值精度}
\end{table}

\subsection{本章小结}

高性能 GPU 编程不仅是关于写内核，更是关于\textbf{正确的开发方法论}：Profile 驱动的优化、理解异步执行、选择合适的维度和精度。这些工程实践建议在日常开发中往往比手写内核更能带来性能收益。

%% ========== Section 11: 总结 ==========

